
\minitoc


% ========================================================================================================
% Introduction 

Jusqu'à maintenant, nous avons étudié des suites de variables aléatoires dites iid (indépendantes et 
identiquement distribuées). Ces propriétés nous ont permis de développer beaucoup d'outils pour les étudier 
tels que la loi forte des grands nombres ou le théorème central limite. 

Or nous n'utilisons pas beaucoup ce type de suite de variables aléatoires en simulation. En effet, ces propriétés 
plutôt fortes ne permettent pas de modéliser des phénomènes dont l'évolution est plus complexe. Nous allons donc 
continuer à étudier des suites de variables aléatoires mais en rajoutant des relations de dépendance entre celles-ci. 
Cela nous permettra de modéliser des systèmes dynnamiques au cours du temps plus complexes. 

% ========================================================================================================
% Définition et Généralités

\section{Définitions et Généralités}

On se place ici dans un espace probabilisé $(E, \mathcal(F), \myP)$  et $E$ un \emph{espace d'état}
fini ou dénombrable. On s'intéresse à une suite de variables aléatoires $(X_n)_{n \in \N}$ ou $X_k$ est 
l'état du système au temps $k$. 

\begin{definition}[Trajectoire]
    On appelle \emph{trajectoire du système} jusqu'au temps $n \in \N^*$ la donnée $(X_1, \dots, X_n)$. 
\end{definition}

\begin{definition}[Chaîne de Markov]
    On dit qu'une suite $(X_n)_{n \geqslant 0}$ est une \emph{chaîne de Markov} si elle 
    vérifie la \emph{propriété de Markov} : 
        \[ \forall n \in \N, \forall i_0, i_1, \dots, i_{n-1} \in E, \] 
        \[ \boxed{ \myP (X_{n+1} = j \; | \; X_n = i, X_{n-1} = i_{n-1}, \dots, X_0 = i_0) = \myP (X_{n+1} = j \; | \; X_n = i)} \]
\end{definition}

\begin{remark}
    Conditionnellement à l'état présent $\{X_n = i\}$, la probabilité d'arriver à l'état $\{X_{n+1} = j\}$ 
    à l'itération suivante ne dépend pas des évènnements du passé $\{X_0 = i_0, \dots, X_{n-1} = n-1\}$. 
    On parle \emph{d'abscence de mémoire}. 
\end{remark}

\begin{definition}[Chaîne de Markov Homogène]
    Soit $(X_n)_{n \geqslant 0}$ une chaîne de Markov. On dit qu'elle est \emph{homogène} si : 
        \[ \forall n \in \N, \forall i,j \in E, \quad \myP (X_{n+1} = j \; | \; X_n = i) = \myP(X_1 = j \; | \; X_0 = i) \]
    On se placera toujours dans ce cadre.
\end{definition}

\begin{remark}
    Une chaîne de Markov homogène est une chaîne de Markov où le temps ($n$) n'apparaît pas dans la probabilité de 
    passer d'un état à l'autre. 
\end{remark}

\begin{example}[Marche Aléatoire]
    On modélise le mouvement d'une particule sur une grille par une chaîne de Markov où $X_n$ représente 
    la position de la particule au temps $n \in  \N$. On considère que la particule se déplace sur une 
    seule dimension. Elle peut donc avancer d'une case ou reculer d'une case à chaque itération. On notera 
    $ p \in [0,1]$ la probabilité qu'elle avance et $1-p$ la probabilité qu'elle recule indépendament du passé. 
    
    Soit $ (\varepsilon_i)_{i \geqslant 0}$ à valeurs dans $\{-1, 1\}$ telle que : 
        \[ \myP( \varepsilon_i = 1) 1 - \myP( \varepsilon_i = -1) = p \in [0,1] \]
    la variable aléatoire qui représente le mouvement de la particule.  
    Déterminons la loi de $X_n$ :
        \[ 
            X_{n+1} = X_n + \varepsilon_{n+1} 
            = X_{n-1} + \varepsilon_n + + \varepsilon_{n+1} 
            = \dots 
            = \sum_{i=1}^{n} \varepsilon_i 
        \] 
    Montrons que c'est bien une chaîne de Markov : 
        \begin{align*}
            X_{n+1} &= X_n + \varepsilon_{n+1} \\ 
            &= \left( \sum_{i=1}^{n} \varepsilon_i \right) + \varepsilon_{n+1} \\ 
            &= \sum_{i=1}^{n+1} \varepsilon_i 
        \end{align*}
    où tous les $ \varepsilon_i, i \in \llbracket 1, n+1 \rrbracket $ sont indépendants. 
    On a donc : 
        \begin{align*}
            \myP(X_{n+1} = j \; | \; X_n = i, \dots, X_0 = i_0) &= \myP(X_n + \varepsilon_{n+1} = j \; | \; X_n = i, \dots, X_0 = i_0) \\ 
            &= \myP(i + \varepsilon_{n+1} = j \; | \; X_n = i, \dots, X_0 = i_0) \\ 
            &= \myP( \varepsilon_{n+1} = j - i) 
        \end{align*}

    On pourra alors s'intéresser au nombre de fois où une trajectoire passe par un point $n \in \N$. 
    La théorie des chaînes de Markov permet de répondre à ce genre de questions. 
\end{example}

Plus généralement, on peut construire récursivement des chaînes de Markov en déinissant un $X_0$ fixé ou tiré par une 
loi $ \pi_0$ puis en posant : 
    \[ X_{n+1} = f(X_n, \varepsilon_{n+1}) \]
où $( \varepsilon_i)$ est une suite de variables aléatoires iid et $f$ une fonction mesurable. 

\begin{definition}[Matrice de Transition]
    Soit $(X_n)_{n \geqslant 0}$ une chaîne de Markov. On définit la \emph{matrice de transition} $P$ 
    de la chaîne comme la matrice, potentiellement infinie, telle que : 
        \[ P(i,j) = \myP(X_1 = j | X_0 = i) \] 
    Ainsi le coefficient de la i-ème ligne et de la j-ème colonne correspond à la probabilité de passer 
    de l'état $i$ à l'état $j$ en une étape.   
\end{definition}

\begin{example}
    Soit $P \in \mathcal{M}_3(\R)$ la matrice de transition suivante : 
        \[ P = 
            \begin{pmatrix}
                1/4 & 3/4 & 0 \\ 
                0 & 1/2 & 1/2 \\ 
                1/2 & 0 & 1/2 
            \end{pmatrix}  \]
    On peut alors la représenter par un graphe pondéré orienté dont les noeuds représentent les 
    états au temps $i$ de la chaîne de Markov et les poids les probabilités de passer d'un état à l'autre en une étape. 
    
    \begin{figure}[h!]
        \centering
        \begin{tikzpicture}[->,>=stealth',auto,node distance=4cm, thick,
                            every loop/.style={min distance=10mm,in=45,out=135,looseness=10}]
            
            % Nodes placés avec des positions relatives
            \node[draw, circle] (1) {1};
            \node[draw, circle, right of=1] (2) {2};
            \node[draw, circle, below of=2] (3) {3};
            
            % Boucles avec labels
            \path (1) edge[loop above] node {$1/4$} (1);
            \path (2) edge[loop right] node {$1/2$} (2);
            \path (3) edge[loop below] node {$1/2$} (3);
            
            % Arêtes avec labels automatiques
            \path (1) edge node {$3/4$} (2);
            \path (2) edge node {$1/2$} (3);
            \path (3) edge node {$1/2$} (1);

        \end{tikzpicture}
        \caption{Graphe Orienté Pondéré de la matrice de transition $P$}
    \end{figure}
\end{example}

\begin{definition}[Matrice Stochastique]
    La matrice de transition $P$ est dite \emph{stochastique} si 
        \[ \forall i \in  E, \quad \sum_{j \in E} P(i,j) = 1 \] 
    (i.e si la somme des poids des arrêtes de tout noeud du graphe associé est égale à 1)
\end{definition}













